% GENERATED FILE. Edit canonical lessons, then run npm run generate:books.
% canonical-source-hash: fnv1a64:ee57ed58648ac977
% canonical-lessons: IT-C379-liceo, IT-C379-frequentare, IT-C379-assenza, IT-C379-pagella, IT-C379-preside

\chapter{A scuola}
\label{ch:it-a-scuola}
\hlchaptermodality{379}

\begin{chapteropening}
\textbf{By the end of this chapter:} \emph{I can talk about secondary school: attending it, absences, the school report and the head teacher.}

Complete the last lesson of chapter 379: I can talk about secondary school: attending it, absences, the school report and the head teacher.
\end{chapteropening}

\section[il liceo]{il liceo — the upper secondary school (academic)}

\label{lesson:IT-C379-liceo}



\begin{quote}
Before the new one: say the Italian for the resolution, then the Italian for the goal.
\end{quote}

\subsection*{You'll want to know: il liceo}
\textbf{il liceo} — \textquotedblleft{}the upper secondary school (academic)\textquotedblright{}.

\begin{quote}
Mia figlia va al liceo scientifico.
\end{quote}

\begin{quote}
\emph{My daughter goes to a science high school.}
\end{quote}

From about fourteen to nineteen. \emph{Il liceo classico} teaches Latin and Greek.

\begin{cousinweb}
\emph{liceo} is a learned borrowing from Latin \emph{lycēum}, from Greek \emph{lýkeion}, the name of a gymnasium with gardens where a great teacher of ancient Greece taught (Lewis \& Short).
\end{cousinweb}

\begin{grammarlens}[title={What you've built}]
One more word for telling, explaining and describing at B1.
\end{grammarlens}

\subsection*{Your turn}
\begin{itemize}
\raggedright
  \item \emph{Say it:} \emph{il liceo}
  \item \emph{Say it:} \emph{il liceo}, once more
  \item \emph{Say it:} a sentence of your own with \emph{il liceo}
  \item \emph{Recall:} say the Italian for the resolution, then the Italian for the goal, then say \emph{il liceo} again
\end{itemize}

\begin{tcolorbox}[breakable,colback=teal!4,colframe=teal!35!black,title={Before you move on}]
What does il liceo mean? (\textquotedblleft{}The upper secondary school (academic)\textquotedblright{}.) Say it once more.
\end{tcolorbox}
\section[frequentare]{frequentare — to attend; to go around with (people)}

\label{lesson:IT-C379-frequentare}



\begin{quote}
Before the new one: say the Italian for the goal, then the Italian for the upper secondary school.
\end{quote}

\subsection*{You'll want to know: frequentare}
\textbf{frequentare} — \textquotedblleft{}to attend; to go around with (people)\textquotedblright{}.

\begin{quote}
Se vuoi passare l'esame, devi frequentare le lezioni.
\end{quote}

\begin{quote}
\emph{If you want to pass the exam, you have to attend the lessons.}
\end{quote}

Also of people: \emph{frequentare brutte compagnie} — to keep bad company.

\begin{cousinweb}
\emph{frequentare} is a learned borrowing from Latin \emph{frequentāre} \textquotedblleft{}to visit often, to crowd\textquotedblright{} (Lewis \& Short), from \emph{frequēns} \textquotedblleft{}crowded, frequent\textquotedblright{}.
\end{cousinweb}

\begin{grammarlens}[title={What you've built}]
One more word for telling, explaining and describing at B1.
\end{grammarlens}

\subsection*{Your turn}
\begin{itemize}
\raggedright
  \item \emph{Say it:} \emph{frequentare}
  \item \emph{Say it:} \emph{frequentare}, once more
  \item \emph{Say it:} a sentence of your own with \emph{frequentare}
  \item \emph{Recall:} say the Italian for the goal, then the Italian for the upper secondary school, then say \emph{frequentare} again
\end{itemize}

\begin{tcolorbox}[breakable,colback=teal!4,colframe=teal!35!black,title={Before you move on}]
What does frequentare mean? (\textquotedblleft{}To attend; to go around with (people)\textquotedblright{}.) Say it once more.
\end{tcolorbox}
\section[l'assenza]{l'assenza — the absence}

\label{lesson:IT-C379-assenza}



\begin{quote}
Before the new one: say the Italian for the upper secondary school, then the Italian for to attend.
\end{quote}

\subsection*{You'll want to know: l'assenza}
\textbf{l'assenza} — \textquotedblleft{}the absence\textquotedblright{}.

\begin{quote}
Se fai troppe assenze, non ti ammettono all'esame.
\end{quote}

\begin{quote}
\emph{If you miss too many days, they will not let you sit the exam.}
\end{quote}

\emph{Fare un'assenza} — to miss a day of school.

\begin{cousinweb}
\emph{assenza} is a learned borrowing from Latin \emph{absentia} \textquotedblleft{}absence\textquotedblright{} (EtymDB), from \emph{absēns} \textquotedblleft{}absent\textquotedblright{}, built on \emph{abesse} \textquotedblleft{}to be away\textquotedblright{}, from \emph{ab} and \emph{esse} \textquotedblleft{}to be\textquotedblright{}.
\end{cousinweb}

\begin{grammarlens}[title={What you've built}]
One more word for telling, explaining and describing at B1.
\end{grammarlens}

\subsection*{Your turn}
\begin{itemize}
\raggedright
  \item \emph{Say it:} \emph{l'assenza}
  \item \emph{Say it:} \emph{l'assenza}, once more
  \item \emph{Say it:} a sentence of your own with \emph{l'assenza}
  \item \emph{Recall:} say the Italian for the upper secondary school, then the Italian for to attend, then say \emph{l'assenza} again
\end{itemize}

\begin{tcolorbox}[breakable,colback=teal!4,colframe=teal!35!black,title={Before you move on}]
What does l'assenza mean? (\textquotedblleft{}The absence\textquotedblright{}.) Say it once more.
\end{tcolorbox}
\section[la pagella]{la pagella — the school report}

\label{lesson:IT-C379-pagella}



\begin{quote}
Before the new one: say the Italian for to attend, then the Italian for the absence.
\end{quote}

\subsection*{You'll want to know: la pagella}
\textbf{la pagella} — \textquotedblleft{}the school report\textquotedblright{}.

\begin{quote}
Ho preso un bel voto in pagella.
\end{quote}

\begin{quote}
\emph{I got a good mark on my school report.}
\end{quote}

Given at the end of each term.

\begin{cousinweb}
\emph{pagella} is a learned borrowing from Latin \emph{pāgella} \textquotedblleft{}a little page\textquotedblright{} (Lewis \& Short), diminutive of \emph{pāgina} \textquotedblleft{}a page\textquotedblright{}.
\end{cousinweb}

\begin{grammarlens}[title={What you've built}]
One more word for telling, explaining and describing at B1.
\end{grammarlens}

\subsection*{Your turn}
\begin{itemize}
\raggedright
  \item \emph{Say it:} \emph{la pagella}
  \item \emph{Say it:} \emph{la pagella}, once more
  \item \emph{Say it:} a sentence of your own with \emph{la pagella}
  \item \emph{Recall:} say the Italian for to attend, then the Italian for the absence, then say \emph{la pagella} again
\end{itemize}

\begin{tcolorbox}[breakable,colback=teal!4,colframe=teal!35!black,title={Before you move on}]
What does la pagella mean? (\textquotedblleft{}The school report\textquotedblright{}.) Say it once more.
\end{tcolorbox}
\section[il preside]{il preside — the head teacher, the principal}

\label{lesson:IT-C379-preside}



\begin{quote}
Before the new one: say the Italian for the absence, then the Italian for the school report.
\end{quote}

\subsection*{You'll want to know: il preside}
\textbf{il preside} — \textquotedblleft{}the head teacher, the principal\textquotedblright{}.

\begin{quote}
Il preside ha chiamato i miei genitori.
\end{quote}

\begin{quote}
\emph{The head teacher called my parents.}
\end{quote}

A woman is \emph{la preside}. The official title today is \emph{dirigente scolastico}.

\begin{cousinweb}
\emph{preside} is a learned borrowing from Latin \emph{praeses, praesidis} \textquotedblleft{}a guardian, a governor\textquotedblright{} (Lewis \& Short), from \emph{praesidēre} \textquotedblleft{}to sit in front\textquotedblright{}, built on \emph{sedēre} \textquotedblleft{}to sit\textquotedblright{}.
\end{cousinweb}

\begin{grammarlens}[title={What you've built}]
One more word for telling, explaining and describing at B1.
\end{grammarlens}

\subsection*{Your turn}
\begin{itemize}
\raggedright
  \item \emph{Say it:} \emph{il preside}
  \item \emph{Say it:} \emph{il preside}, once more
  \item \emph{Say it:} a sentence of your own with \emph{il preside}
  \item \emph{Recall:} say the Italian for the absence, then the Italian for the school report, then say \emph{il preside} again
\end{itemize}

\begin{tcolorbox}[breakable,colback=teal!4,colframe=teal!35!black,title={Before you move on}]
What does il preside mean? (\textquotedblleft{}The head teacher, the principal\textquotedblright{}.) Say it once more.
\end{tcolorbox}
